% =====================================================================
\section{Знакомство с ESP32-S3 и настройка среды}
% =====================================================================
\lessonintro
  {понять, из чего состоит ESP32-S3, как устроена отладочная плата, установить среду и залить первую программу}
  {ничего --- это стартовый урок}
  {знакомимся с проектом «умная лампа», который будем собирать весь курс}
  {плата ESP32-S3-DevKitC-1, USB-кабель с передачей данных}

\subsection{Что такое ESP32-S3}

ESP32-S3 --- система на кристалле (SoC) компании Espressif. Главное:

\begin{itemize}
  \item \textbf{Два ядра Xtensa LX7} до \SI{240}{MHz} с векторными инструкциями (ускорение DSP и нейросетей).
  \item \textbf{\SI{512}{KB} SRAM} на кристалле, \SI{384}{KB} ROM, до \SI{16}{MB} внешней Flash и
        до \SI{16}{MB} внешней PSRAM (в том числе в режиме Octal SPI).
  \item \textbf{Wi-Fi 802.11 b/g/n} (\SI{2.4}{GHz}) и \textbf{Bluetooth 5 LE}. Классического Bluetooth (BR/EDR)
        \emph{нет} --- в отличие от первого ESP32.
  \item \textbf{Нативный USB OTG} и встроенный \textbf{USB-Serial/JTAG}: прошивать и отлаживать можно
        без отдельного преобразователя.
  \item \textbf{45 GPIO}, 2 АЦП по 12 бит (20 каналов), 14 сенсорных (touch) входов, 3 UART, 2 I\textsuperscript{2}C,
        2 I\textsuperscript{2}S, 4 SPI (два доступны пользователю), LEDC, MCPWM, RMT, TWAI (CAN), интерфейс камеры и LCD.
  \item \textbf{Сопроцессоры ULP} (RISC-V и FSM), которые работают, пока основные ядра спят.
\end{itemize}

Не нужно запоминать всё сразу: каждый блок схемы ниже мы по очереди задействуем в уроках курса.

\begin{figure}[H]
\centering
\begin{tikzpicture}[node distance=3mm,
  box/.style={blk,minimum width=26mm}]
  \node[box,fill=blkcpu,minimum width=30mm,minimum height=11mm] (c0) {Xtensa LX7\\ ядро 0 (PRO)};
  \node[box,fill=blkcpu,minimum width=30mm,minimum height=11mm,right=4mm of c0] (c1) {Xtensa LX7\\ ядро 1 (APP)};
  \node[box,fill=blkmem,below=6mm of c0,minimum width=30mm] (sram) {SRAM 512 КБ};
  \node[box,fill=blkmem,below=6mm of c1,minimum width=30mm] (rom) {ROM 384 КБ};
  \node[box,fill=blkmem,below=3mm of sram,minimum width=30mm] (cache) {Кэш};
  \node[box,fill=blkmem,below=3mm of rom,minimum width=30mm] (mspi) {Контроллер\\ Flash / PSRAM};
  \coordinate (busL) at ($(cache.south west)+(-2.4,-0.6)$);
  \coordinate (busR) at ($(mspi.south east)+(2.4,-0.6)$);
  \draw[line width=4pt,gray!45] (busL) -- (busR) node[midway,below=1pt,font=\sffamily\scriptsize,black] {внутренняя шина};
  \node[box,fill=blkrf,left=10mm of c0,minimum width=24mm,minimum height=11mm] (wifi) {Wi-Fi\\ 802.11 b/g/n};
  \node[box,fill=blkrf,below=3mm of wifi,minimum width=24mm,minimum height=11mm] (bt) {Bluetooth\\ 5 LE};
  \node[box,fill=blkrf,below=3mm of bt,minimum width=24mm] (rf) {RF + балун};
  \node[box,fill=blkrtc,right=10mm of c1,minimum width=24mm,minimum height=11mm] (ulp) {ULP-RISC-V\\ ULP-FSM};
  \node[box,fill=blkrtc,below=3mm of ulp,minimum width=24mm] (rtcm) {RTC-память\\ 16 КБ};
  \node[box,fill=blkrtc,below=3mm of rtcm,minimum width=24mm] (pmu) {Управление\\ питанием};
  \node[box,fill=blkper,minimum width=18mm,anchor=north west] (p1) at ($(busL)+(0.2,-0.6)$) {GPIO\\ ×45};
  \node[box,fill=blkper,minimum width=18mm,right=2mm of p1] (p2) {UART\\ ×3};
  \node[box,fill=blkper,minimum width=18mm,right=2mm of p2] (p3) {I\textsuperscript{2}C ×2\\ I\textsuperscript{2}S ×2};
  \node[box,fill=blkper,minimum width=18mm,right=2mm of p3] (p4) {SPI\\ ×4};
  \node[box,fill=blkper,minimum width=18mm,right=2mm of p4] (p5) {ADC ×2\\ Touch ×14};
  \node[box,fill=blkper,minimum width=18mm,right=2mm of p5] (p6) {LEDC\\ MCPWM};
  \node[box,fill=blkper,minimum width=18mm,right=2mm of p6] (p7) {USB OTG\\ USB-JTAG};
  \node[box,fill=blkper,minimum width=18mm,right=2mm of p7] (p8) {RMT, TWAI\\ LCD/CAM};
  \draw[darr] (c0) -- (sram);
  \draw[darr] (c1) -- (rom);
  \draw[darr] (c0.east) -- (c1.west);
  \draw[darr] (wifi.east) -- (c0.west);
  \foreach \p in {p1,p2,p3,p4,p5,p6,p7,p8} \draw[thick,gray] (\p.north) -- (\p.north |- busL);
  \draw[thick,gray] (cache.south) -- (cache.south |- busL);
  \draw[thick,gray] (mspi.south) -- (mspi.south |- busL);
  \draw[darr] (ulp.west) -- (c1.east);
  \begin{scope}[on background layer]
    \node[draw=espdark,thick,rounded corners=4pt,fit=(wifi)(ulp)(p1)(p8)(pmu)(rf),inner sep=5pt,
      label={[font=\sffamily\bfseries,text=espdark]above:ESP32-S3}] {};
  \end{scope}
\end{tikzpicture}
\caption{Упрощённая блок-схема ESP32-S3}
\end{figure}

\subsection{Сравнение с «родственниками»}

\begin{table}[H]
\centering\small
\begin{tabularx}{\textwidth}{@{}l X X X@{}}
\toprule
 & \textbf{ESP32} & \textbf{ESP32-S3} & \textbf{ESP32-C3} \\
\midrule
Ядро & 2× Xtensa LX6, 240 МГц & 2× Xtensa LX7, 240 МГц & 1× RISC-V, 160 МГц \\
SRAM & 520 КБ & 512 КБ & 400 КБ \\
Bluetooth & Classic + BLE 4.2 & BLE 5 & BLE 5 \\
USB & нет & OTG + Serial/JTAG & Serial/JTAG \\
GPIO & 34 & 45 & 22 \\
Особенности & ЦАП, Ethernet MAC & векторные инструкции, камера, LCD & дешёвый, малый \\
\bottomrule
\end{tabularx}
\caption{Сравнение популярных чипов Espressif}
\end{table}

\subsection{Модули и маркировка}

Чип почти всегда используется в составе \emph{модуля} (ESP32-S3-WROOM-1, WROOM-1U, MINI-1),
где уже есть Flash, кварц и антенна. Суффикс модуля кодирует объём памяти:

\begin{center}
\begin{tikzpicture}[font=\sffamily]
  \node[font=\sffamily\Large\bfseries] (m) {ESP32-S3-WROOM-1-\textcolor{accent}{N16}\textcolor{esp}{R8}};
  \draw[accent,thick,decorate,decoration={brace,mirror,amplitude=4pt}] ($(m.south east)+(-1.25,-0.05)$) -- ($(m.south east)+(-0.62,-0.05)$)
     node[midway,below=5pt,align=center,font=\sffamily\small,text=accent] {Flash 16 МБ};
  \draw[esp,thick,decorate,decoration={brace,mirror,amplitude=4pt}] ($(m.south east)+(-0.58,-0.05)$) -- ($(m.south east)+(0,-0.05)$)
     node[midway,below=5pt,xshift=12mm,align=center,font=\sffamily\small,text=esp] {PSRAM 8 МБ\\ (Octal)};
\end{tikzpicture}
\end{center}

\begin{caution}
В модулях с \textbf{Octal PSRAM} (\code{R8}, \code{R16V}) выводы \pin{35}, \pin{36}, \pin{37} заняты
памятью --- не используйте их. Выводы \pin{26}--\pin{32} всегда заняты SPI-Flash и на плату обычно не выведены.
\end{caution}

\subsection{Отладочная плата ESP32-S3-DevKitC-1}

\begin{figure}[H]
\centering
\begin{tikzpicture}[x=1mm,y=1mm,font=\sffamily\scriptsize]
  \def\pitch{4.2}
  \fill[black!80,rounded corners=2mm] (0,0) rectangle (46,98);
  \fill[gray!55,rounded corners=0.6mm] (7,50) rectangle (39,94);
  \fill[gray!35] (7,84) rectangle (39,94);
  \node[white,font=\sffamily\bfseries\scriptsize] at (23,68) {ESP32-S3};
  \node[white,font=\sffamily\tiny] at (23,64) {WROOM-1};
  \node[black,font=\sffamily\tiny] at (23,89) {антенна};
  \fill[gray!70] (8,-3) rectangle (18,5);  \node[white] at (13,8) {UART};
  \fill[gray!70] (28,-3) rectangle (38,5); \node[white] at (33,8) {USB};
  \fill[gray!40] (4,14) rectangle (10,20); \node[white] at (7,23) {BOOT};
  \fill[gray!40] (36,14) rectangle (42,20); \node[white] at (39,23) {RST};
  \fill[white] (21,30) rectangle (25,34);
  \node[white] at (23,27) {RGB};
  \foreach \lab [count=\i] in {3V3,3V3,RST,4,5,6,7,15,16,17,18,8,3,46,9,10,11,12,13,14,5V,GND}{
    \pgfmathsetmacro{\yy}{96-(\i-1)*\pitch}
    \fill[yellow!70!orange] (2.2,\yy) circle (1);
    \node[anchor=east] (L\i) at (-2,\yy) {\lab};
  }
  \foreach \lab [count=\i] in {GND,43 TX,44 RX,1,2,42,41,40,39,38,37,36,35,0,45,48,47,21,20,19,GND,GND}{
    \pgfmathsetmacro{\yy}{96-(\i-1)*\pitch}
    \fill[yellow!70!orange] (43.8,\yy) circle (1);
    \node[anchor=west] (R\i) at (48,\yy) {\lab};
  }
  \node[font=\sffamily\bfseries\small,anchor=south] at (-9,99) {J1};
  \node[font=\sffamily\bfseries\small,anchor=south] at (55,99) {J3};
  \foreach \i in {1,2,21}{\node[draw=esp,thick,rounded corners=1pt,fit=(L\i),inner sep=0.3pt]{};}
  \foreach \i in {22}{\node[draw=gray,thick,rounded corners=1pt,fit=(L\i),inner sep=0.3pt]{};}
  \foreach \i in {1,21,22}{\node[draw=gray,thick,rounded corners=1pt,fit=(R\i),inner sep=0.3pt]{};}
  \foreach \i in {14,15}{\node[draw=warn,thick,rounded corners=1pt,fit=(R\i),inner sep=0.3pt]{};}
  \foreach \i in {13,14}{\node[draw=warn,thick,rounded corners=1pt,fit=(L\i),inner sep=0.3pt]{};}
  \foreach \i in {11,12,13}{\node[draw=blkrtc!40!black,thick,dashed,rounded corners=1pt,fit=(R\i),inner sep=0.3pt]{};}
  \foreach \i in {19,20}{\node[draw=accent,thick,rounded corners=1pt,fit=(R\i),inner sep=0.3pt]{};}
  % выводы проекта
  \foreach \i in {4,5,12}{\node[fill=okgreen,circle,inner sep=1.3pt] at ($(L\i.west)+(-1.5,0)$) {};}
  \foreach \i in {4}{\node[fill=okgreen,circle,inner sep=1.3pt] at ($(R\i.east)+(1.5,0)$) {};}
  \node[fill=okgreen,circle,inner sep=1.3pt] at ($(L15.west)+(-1.5,0)$) {};
  \node[anchor=west,align=left,font=\sffamily\scriptsize] at (68,80) {%
    \textcolor{esp}{\rule{3mm}{2mm}} питание\\[1pt]
    \textcolor{gray}{\rule{3mm}{2mm}} земля\\[1pt]
    \textcolor{warn}{\rule{3mm}{2mm}} strapping-выводы\\[1pt]
    \textcolor{accent}{\rule{3mm}{2mm}} USB D$-$/D$+$ (19/20)\\[1pt]
    \textcolor{blkrtc!40!black}{\rule{3mm}{2mm}} заняты Octal PSRAM\\[1pt]
    \textcolor{okgreen}{$\bullet$} выводы проекта «умная лампа»\\[6pt]
    RGB-светодиод WS2812:\\ GPIO48 (v1.0) или GPIO38 (v1.1)\\[4pt]
    Кнопка BOOT = GPIO0\\[4pt]
    ADC1: GPIO1--10\\ ADC2: GPIO11--20\\[4pt]
    Уровни только \textbf{3,3 В}!};
\end{tikzpicture}
\caption{Разводка выводов ESP32-S3-DevKitC-1 (вид сверху). Числа --- номера GPIO.
Сверьтесь с документацией именно вашей ревизии платы.}
\end{figure}

У платы \textbf{два USB-разъёма}:
\begin{itemize}
  \item \textbf{UART} --- через микросхему-преобразователь (CP2102N/CH343) на UART0 (\pin{43}/\pin{44}).
        Самый надёжный вариант для прошивки и вывода сообщений.
  \item \textbf{USB} --- напрямую к выводам нативного USB чипа (\pin{19}/\pin{20}). Позволяет
        прошивать, отлаживать через JTAG и превращать плату в USB-клавиатуру, MIDI-устройство и т.\,п.
\end{itemize}

\subsection{Установка среды}

Для уроков используем Arduino IDE 2.x (подойдёт и PlatformIO --- код не меняется).

\begin{enumerate}
  \item Установите Arduino IDE 2 с сайта \url{https://www.arduino.cc/en/software}.
  \item \textsf{File → Preferences → Additional boards manager URLs}, добавьте:\\
    {\small\url{https://espressif.github.io/arduino-esp32/package_esp32_index.json}}
  \item \textsf{Tools → Board → Boards Manager}, найдите \textbf{esp32 by Espressif Systems},
        установите версию 3.x.
  \item Выберите плату \textbf{ESP32S3 Dev Module} и порт.
  \item Если подключены через разъём \textbf{USB} (нативный), включите \textsf{Tools → USB CDC On Boot: Enabled},
        иначе сообщения от платы не дойдут до компьютера.
  \item Для модулей с PSRAM выберите \textsf{Tools → PSRAM: OPI PSRAM} (для \code{R8}) или \textsf{QSPI PSRAM} (для \code{R2}).
\end{enumerate}

\begin{figure}[H]
\centering
\begin{tikzpicture}[node distance=6mm and 8mm]
  \node[blk,fill=blkmem] (src) {Скетч\\ \code{.ino}};
  \node[blk,fill=blkmem,right=of src] (cc) {Компилятор\\ xtensa-esp32s3-gcc};
  \node[blk,fill=blkmem,right=of cc] (bin) {Образ\\ \code{.bin}};
  \node[blk,fill=blkper,right=of bin] (tool) {esptool};
  \node[blk,fill=blkcpu,right=of tool] (chip) {Flash\\ ESP32-S3};
  \draw[arr] (src) -- (cc);
  \draw[arr] (cc) -- (bin);
  \draw[arr] (bin) -- (tool);
  \draw[arr] (tool) -- node[above,font=\sffamily\scriptsize]{USB} (chip);
  \node[blk,fill=blkrf,below=8mm of cc] (lib) {ESP-IDF + FreeRTOS\\ Arduino-ядро};
  \draw[arr] (lib) -- (cc);
\end{tikzpicture}
\caption{Путь от исходного кода до прошивки}
\end{figure}

\subsubsection{Если прошивка не начинается}

Чип входит в режим загрузчика, если при сбросе \pin{0} притянут к земле. Обычно IDE делает это
автоматически, но если видите \code{Failed to connect}:

\begin{center}
\begin{tikzpicture}[node distance=4mm]
  \node[blk,fill=blkper] (a) {Зажать\\ BOOT};
  \node[blk,fill=blkper,right=of a] (b) {Нажать и отпустить\\ RST};
  \node[blk,fill=blkper,right=of b] (c) {Отпустить\\ BOOT};
  \node[blk,fill=blkrf,right=of c] (d) {Нажать\\ «Upload»};
  \node[blk,fill=blkcpu,right=of d] (e) {После прошивки\\ нажать RST};
  \draw[arr] (a)--(b); \draw[arr] (b)--(c); \draw[arr] (c)--(d); \draw[arr] (d)--(e);
\end{tikzpicture}
\end{center}

\subsection{Структура Arduino-скетча}

Любой скетч состоит из двух функций: \code{setup()} выполняется один раз после включения,
\code{loop()} --- повторяется бесконечно. (На самом деле под ними работает операционная система
FreeRTOS; она нам понадобится только в уроке~7.)

\begin{center}
\begin{tikzpicture}[node distance=6mm]
  \node[blk,fill=blkrtc] (rst) {Сброс / подача питания};
  \node[blk,fill=blkmem,below=of rst] (boot) {ROM-загрузчик → 2-й загрузчик → ESP-IDF};
  \node[blk,fill=blkper,below=of boot] (setup) {\code{setup()} --- один раз};
  \node[blk,fill=blkper,below=of setup] (loop) {\code{loop()} --- бесконечно};
  \draw[arr] (rst)--(boot); \draw[arr] (boot)--(setup); \draw[arr] (setup)--(loop);
  \draw[arr] (loop.east) to[out=0,in=0,looseness=2.5] (loop.north east);
\end{tikzpicture}
\end{center}

Первая программа ничего не подключает: она только сообщает о себе через USB. Как именно работает
этот канал связи, разберём в уроке~2, а пока просто проверим, что всё настроено.

\codecap{Первый скетч: «Hello, ESP32-S3»}
\codeinput{l0_hello}

Откройте \textsf{Tools → Serial Monitor}, установите скорость 115200 --- должны появиться сообщения.

\subsection{Сквозной проект: умная лампа}
\label{sec:project}

Чтобы уроки не были набором разрозненных примеров, все они работают на один проект. К концу курса
у нас будет настольная лампа, которая:

\begin{itemize}
  \item включается кнопкой и командой с компьютера;
  \item плавно меняет яркость ручкой-потенциометром;
  \item показывает своё состояние на OLED-экране;
  \item управляется с телефона через браузер;
  \item засыпает, когда выключена, и просыпается от кнопки.
\end{itemize}

Каждый урок добавляет ровно одну возможность, опираясь на уже изученное. Выводы платы закреплены
за деталями с самого начала и дальше не меняются:

\begin{figure}[H]
\centering
\begin{circuitikz}[scale=0.95,transform shape]
  \node[draw,thick,fill=blkcpu,minimum width=3.2cm,minimum height=6.2cm,align=center,font=\sffamily] (esp) at (0,0)
     {ESP32-S3\\ DevKitC-1\\[4mm] \scriptsize RGB: GPIO48\\ \scriptsize (урок 1)};
  % LED
  \draw ($(esp.east)+(0,2.3)$) node[left,font=\sffamily\scriptsize]{GPIO4} to[short,-] (1.9,2.3)
        to[R, l=\SI{220}{\ohm}] (4.2,2.3) to[led, l=LED лампы] (6.4,2.3) -- (7,2.3) node[right,font=\sffamily\small]{GND};
  \node[font=\sffamily\scriptsize,text=okgreen] at (9.2,2.3) {уроки 1, 4};
  % button
  \draw ($(esp.east)+(0,1.1)$) node[left,font=\sffamily\scriptsize]{GPIO5} to[short,-] (4.2,1.1)
        to[push button, l=S1] (6.4,1.1) -- (7,1.1) node[right,font=\sffamily\small]{GND};
  \node[font=\sffamily\scriptsize,text=okgreen] at (9.2,1.1) {урок 3};
  % pot
  \draw (3.4,-0.9) node[left,font=\sffamily\small]{3,3 В} to[pR, n=pot, l_=\SI{10}{k\ohm}] (6.4,-0.9) -- (7,-0.9) node[right,font=\sffamily\small]{GND};
  \draw ($(esp.east)+(0,-0.1)$) node[left,font=\sffamily\scriptsize]{GPIO1} -| (pot.wiper);
  \node[font=\sffamily\scriptsize,text=okgreen] at (9.2,-0.9) {урок 5};
  % OLED
  \node[draw,thick,fill=blkper,minimum width=2.2cm,minimum height=1.5cm,align=center,font=\sffamily\small] (oled) at (6,-2.35) {OLED\\ SSD1306};
  \draw ($(esp.east)+(0,-2.0)$) node[left,font=\sffamily\scriptsize]{GPIO8} -- node[above,font=\sffamily\scriptsize,pos=0.3]{SDA} (oled.west |- 0,-2.0);
  \draw ($(esp.east)+(0,-2.7)$) node[left,font=\sffamily\scriptsize]{GPIO9} -- node[below,font=\sffamily\scriptsize,pos=0.3]{SCL} (oled.west |- 0,-2.7);
  \node[font=\sffamily\scriptsize,text=okgreen] at (9.2,-2.35) {урок 6};
  \node[font=\sffamily\scriptsize,align=left,anchor=west] at (7.4,-3.3) {питание экрана: 3V3, GND};
\end{circuitikz}
\caption{Итоговая схема умной лампы. Детали добавляются по мере прохождения уроков}
\end{figure}

\begin{figure}[H]
\centering
\begin{tikzpicture}[node distance=4mm and 3mm,
  st/.style={blk,fill=blkrf,minimum width=27mm,minimum height=13mm,font=\sffamily\scriptsize}]
  \node[st] (s1) {\textbf{Урок 1}\\ светодиод лампы,\\ «пульс» на RGB};
  \node[st,right=of s1] (s2) {\textbf{Урок 2}\\ команды on/off\\ с компьютера};
  \node[st,right=of s2] (s3) {\textbf{Урок 3}\\ кнопка-\\выключатель};
  \node[st,right=of s3] (s4) {\textbf{Урок 4}\\ плавная яркость\\ (ШИМ)};
  \node[st,right=of s4] (s5) {\textbf{Урок 5}\\ ручка яркости\\ (АЦП)};
  \node[st,below=9mm of s5] (s6) {\textbf{Урок 6}\\ экран состояния\\ (I\textsuperscript{2}C)};
  \node[st,left=of s6] (s7) {\textbf{Урок 7}\\ разделение\\ на задачи};
  \node[st,left=of s7] (s8) {\textbf{Урок 8}\\ управление\\ через браузер};
  \node[st,left=of s8] (s9) {\textbf{Урок 9}\\ сон и пробуждение\\ кнопкой};
  \node[st,left=of s9,fill=blkcpu] (fin) {\textbf{Приложение}\\ полный код\\ проекта};
  \draw[arr] (s1)--(s2); \draw[arr] (s2)--(s3); \draw[arr] (s3)--(s4); \draw[arr] (s4)--(s5);
  \draw[arr] (s5)--(s6); \draw[arr] (s6)--(s7); \draw[arr] (s7)--(s8); \draw[arr] (s8)--(s9); \draw[arr] (s9)--(fin);
\end{tikzpicture}
\caption{Как растёт проект от урока к уроку}
\end{figure}

\begin{summary}
\begin{itemize}[leftmargin=*]
  \item ESP32-S3 --- двухъядерный чип с Wi-Fi, BLE, USB и богатой периферией; работает на \SI{3.3}{V}.
  \item Прошивка заливается из Arduino IDE; при проблемах помогает ручной вход в загрузчик (BOOT + RST).
  \item Скетч --- это \code{setup()} (один раз) и \code{loop()} (бесконечно).
\end{itemize}
\end{summary}

\begin{tasks}
\begin{enumerate}
  \item Выясните, какой модуль стоит на вашей плате (надпись на экране модуля) и сколько у него Flash и PSRAM. Сверьте с выводом скетча.
  \item Выведите MAC-адрес чипа: \code{ESP.getEfuseMac()}.
  \item Измените частоту CPU функцией \code{setCpuFrequencyMhz(80)} и проверьте результат.
\end{enumerate}
\end{tasks}
